% ARQUIVO GERADO AUTOMATICAMENTE por article/scripts/gerar_resultados.py
% Fonte: data/outputs/*.parquet — NÃO EDITAR À MÃO.
\begin{table}[htbp]
\centering
\caption{Valores ausentes no questionário e robustez da imputação pela moda}
\label{tab:qualidade_dados}
\resizebox{\textwidth}{!}{%
\begin{tabular}{ccccccc}
\toprule
\textbf{Edição} & \textbf{Participantes} & \textbf{Questionário nulo} & \textbf{\% nulos} & \textbf{Teste / base} & \textbf{Menor margem da moda} & \textbf{Moda robusta} \\
\midrule
2019 & 3.701.910 & 0 & 0,00\% & 0,27\% & \texttt{Q006} (12.668) & Sim \\
2020 & 2.588.681 & 27.377 & 1,06\% & 0,39\% & \texttt{Q024} (47.606) & Sim \\
2021 & 2.238.107 & 1 & 0,00\% & 0,45\% & \texttt{TP\_COR\_RACA} (66.595) & Sim \\
\bottomrule
\end{tabular}}
\fonte{Elaborado pelo autor a partir de \texttt{data\_quality.parquet}. Questionário nulo: registros com Q001, Q002, Q006, Q024 e Q025 simultaneamente ausentes. Teste / base: fração que o conjunto de teste (10.000) representa na população usada para calcular a moda. Menor margem: diferença, em número de registros, entre a moda e a segunda categoria mais frequente na coluna mais apertada. Moda robusta: a margem supera 10.000 em todas as colunas, de modo que o teste não poderia alterar a moda.}
\end{table}
